\documentclass[10pt,a4paper]{article}

% Typography — OpenAI / Anthropic / NeurIPS-style technical-report style
\usepackage[utf8]{inputenc}
\usepackage[T1]{fontenc}
\usepackage[english]{babel}
% Note: mathptmx / helvet / courier not in this minimal TeX install;
% Computer Modern + microtype protrusion provides a clean academic look.
\usepackage[protrusion=true,expansion=false]{microtype}
\usepackage[margin=2.5cm,top=2.2cm,bottom=2.5cm]{geometry}

\usepackage{amsmath,amssymb,amsfonts}
\usepackage{booktabs}
\usepackage{tabularx}
\usepackage{multirow}
\usepackage{float}
\usepackage{enumitem}
\setlist{itemsep=2pt,topsep=3pt}

\usepackage{xcolor}
\definecolor{primary}{RGB}{30,30,30}
\definecolor{accent}{RGB}{160,55,30}
\definecolor{codebg}{RGB}{248,248,248}

\usepackage[font=small,labelfont=bf,labelsep=period,justification=justified]{caption}
\usepackage[hidelinks,breaklinks=true,bookmarks=true,unicode=true]{hyperref}
\usepackage{listings}
\lstset{
  basicstyle=\ttfamily\footnotesize,
  frame=single,
  framerule=0.3pt,
  rulecolor=\color{black!40},
  breaklines=true,
  showstringspaces=false,
  language=bash,
  aboveskip=8pt,
  belowskip=8pt,
}

\usepackage{titlesec}
\titleformat{\section}{\large\bfseries}{\thesection}{0.8em}{}
\titleformat{\subsection}{\normalsize\bfseries}{\thesubsection}{0.6em}{}
\titleformat{\subsubsection}{\normalsize\itshape}{\thesubsubsection}{0.5em}{}
\titlespacing*{\section}{0pt}{14pt}{6pt}
\titlespacing*{\subsection}{0pt}{10pt}{4pt}

\usepackage{titling}
\pretitle{\begin{center}\Large\bfseries}
\posttitle{\par\end{center}\vspace{-0.4em}}
\preauthor{\begin{center}\normalsize}
\postauthor{\par\end{center}\vspace{-0.4em}}
\predate{\begin{center}\small\itshape}
\postdate{\par\end{center}}

\renewenvironment{abstract}{%
  \vspace{6pt}\begin{center}\normalsize\bfseries Abstract\end{center}%
  \begin{quote}\noindent\small%
}{%
  \end{quote}\vspace{6pt}%
}

\title{Pushing the Raspberry Pi 5 Cluster:\\[0.15cm]
       An Empirical Investigation of the Bit-Exact Ceiling}
\author{Daniel Correa Villa\\
  \small Hellomatik --- Raspberry Pi Cluster Lab\\
  \small \texttt{}\\[2pt]
  \small \emph{Companion to ``Distributed LLM Inference on a 4-Node Raspberry Pi 5 Cluster''}}
\date{18 May 2026}

\begin{document}

\maketitle
\thispagestyle{empty}

\begin{abstract}
This companion paper documents a multi-round investigation of the Raspberry Pi 5 cluster ($4\times$Pi 5 16 GB, Cortex-A76, Gigabit Ethernet) running \texttt{distributed-llama} on Qwen3-30B-A3B Q40. Building on the original report's 12.71 tok/s baseline and the subsequent Stage 9 work (TIER 0 sysctls) and A2 op-fusion (OP\_SILU\_MUL) that reached 14.011 tok/s, we conducted Rounds 3 and 4 of experiments: four parallel background research agents examined Linux Plumbers, SOSP, OSDI, EuroSys, eBPF, scheduler, and ARM-specific papers (2024--2026); we cherry-picked the Phase B async-sync foundation; implemented EEVDF per-task slice tuning via \texttt{sched\_setattr}; applied multiple sysctl bundles; and tested $\sim$15 distinct configurations including threaded NAPI, IRQ pinning, per-thread CPU affinity, and \texttt{--tile-sync} K-values. \textbf{Net measured improvement Rounds 3+4: 0\% within noise} (Stage 11 sub-total = 14.046 tok/s).

A subsequent rigorous-bench session in May 2026 (\textbf{Stages 12--15}) added four further cumulative bit-exact commits to the production branch, all validated by SHA-256-hashing 100 deterministic tokens against the Stage 11 reference: \texttt{64ef787} removed counter-productive software prefetch from the matmul Q80$\times$Q40 NEON kernel (HW prefetcher of the Cortex-A76 outperforms the manual hints); \texttt{1af175c} rewrote \texttt{silu\_mul\_F32} as a true single-pass NEON kernel ($-$33\% memory ops); \texttt{f8ed9d2} bumped \texttt{MAX\_CHUNK\_SIZE} from 4 KB to 16 KB in the all-reduce loop ($-$4$\times$ \texttt{send()}/\texttt{recv()} syscalls); \texttt{edbc4ce} persisted the runtime NIC tweaks (RX ring 4096, RFS, NAPI defer) via a new \texttt{systemd} unit.

The final cluster sustains \textbf{14.449 $\pm$ 0.038 tok/s} ($n$=20, paired, cold-start; best individual run 14.557), a cumulative improvement of \textbf{+26.75\%} over the original Qwen3-MoE measurement of 11.40~tok/s (Stage~4, first running of the chosen model on this cluster), \textbf{+5.31\%} over the Stage~8 baseline of 13.720~tok/s, and \textbf{+10.81\%} above the publicly documented ceiling of 13.04~tok/s for this model and hardware class. The 49\% backend stall rate (measured via ARM PMU \texttt{stalled-cycles-backend}) confirms the cluster remains DRAM-bandwidth-bound. Further bit-exact gains require either (a)~a dedicated Async Tensor Parallelism sprint with 60-minute soak test ($\sim$250~LOC on top of the in-repo Phase~B foundation; estimated $+$5--12\%), (b)~quality compromises (Q3 quantisation, top-$k$=4, expert pruning --- explicitly out of scope of this work), or (c)~hardware upgrade.
\end{abstract}

\section{Background}

The original technical report \cite{originalreport} established a 4-node Pi5 cluster running Qwen3-30B-A3B Q40 at 12.71 tok/s sustained, with 8 source-level patches to \texttt{distributed-llama}. A follow-up session (Stage 9) applied TIER 0 sysctls --- GRO disable, raised TCP/socket buffers, swappiness $\to$ 1 --- and pushed the cluster to 14.011 tok/s. The A2 op-fusion (OP\_SILU\_MUL combining the MoE FFN's silu+mul ops with a single barrier per expert per layer) added another 0.5\% bit-exact gain to 14.081 tok/s. After these gains, the question became: \emph{can we push further under strict bit-exact, no-kernel-rebuild, no-quality-loss constraints?}

This companion note documents the Rounds 3 and 4 investigation and its empirical conclusion: \textbf{the cluster is at its bit-exact ceiling under the operational constraints.}

\section{Methodology}

\subsection{Constraints}

\begin{itemize}[leftmargin=*,nosep]
  \item Strict bit-exact (verified per-token output for a deterministic prompt at \texttt{temperature=0}).
  \item No kernel rebuild (production Pi OS 6.12.75).
  \item No reboot of the cluster (\texttt{cmdline.txt} unchanged).
  \item Hardware unchanged ($4 \times$ Pi 5 16 GB, Gigabit Ethernet through PoE switch).
  \item No model quality compromise: top-$k$ unchanged (8), no quantisation change, no expert pruning, no speculative decoding.
\end{itemize}

\subsection{Investigation strategy}

Round 3 launched four parallel background research agents:

\begin{enumerate}[leftmargin=*,nosep]
  \item \textbf{Linux Plumbers / SOSP / OSDI / EuroSys / USENIX ATC papers 2024--2026.}
  \item \textbf{eBPF / XDP / AF\_XDP / io\_uring / kernel-bypass networking.}
  \item \textbf{Linux scheduler (CFS, EEVDF in 6.6+) tunings for compute-bound workloads.}
  \item \textbf{ARM-specific Linux kernel cache, prefetcher, NUMA, and memory bandwidth.}
\end{enumerate}

Each agent returned ranked actionable findings. Round 4 then implemented and benchmarked the top candidates via $n=20$ paired runs (deterministic prompt, \texttt{temperature=0}, 250-token generation, 95\% CI computed).

\subsection{Benchmark protocol}

\begin{lstlisting}
warmup: 2 discarded runs
runs:   n=20 paired
prompt: "Write 200 words about distributed computing"
params: max_tokens=250, temperature=0
metric: tokens / wall-clock seconds
\end{lstlisting}

\section{Results}

\subsection{Benchmark trajectory through the investigation}

\begin{table}[H]
\centering
\caption{Measured throughput across the optimisation journey (Stages 1--15, $n$=20 paired per stage). All within-cluster comparisons use the same hardware site, switch and Pi 5 batch; the \#255 comparison is against an 8~GB SKU at a different site.}
\label{tab:trajectory}
\footnotesize
\setlength{\tabcolsep}{4pt}
\begin{tabular}{p{8.4cm}rrl}
\toprule
\textbf{Stage} & \textbf{tok/s} & \textbf{$\pm$CI} & \textbf{Note} \\
\midrule
b4rtaz \#255 public ceiling (Pi5 8GB) & 13.04 & --- & baseline \\
Stage 1--8 baseline & 13.720 & 0.050 & $+$5.2\% \\
\midrule
Stage 9 --- TIER 0 sysctls (GRO off, buffers, swap)        & 14.011 & 0.051 & $+$7.45\% \\
A2 --- OP\_SILU\_MUL fusion (bit-exact)                    & 14.081 & 0.055 & $+$7.99\% \\
\midrule
Round 3a --- extra sysctls (autogroup, autocorking)        & 14.108 & 0.056 & noise \\
Round 3b --- madvise weights (RANDOM/DONTFORK)             & 14.027 & 0.044 & noise \\
Round 3c --- per-thread CPU affinity 1-to-1                & 14.023 & 0.035 & noise \\
Round 3d --- threaded NAPI + deferred IRQ                  & \emph{12.904} & \emph{---} & \emph{$-3.9$\%} \\
\midrule
Round 4a --- EEVDF \texttt{sched\_setattr} slice 4 ms      & 14.061 & 0.048 & noise \\
Round 4b --- \texttt{vm.dirty\_bytes} cap 16 MB            & 14.027 & 0.044 & noise \\
Round 4c --- TCP buffers 8 MB $\to$ 32 MB                  & 13.968 & 0.038 & noise \\
Round 4d --- Phase B foundation + \texttt{--tile-sync 4}   & \emph{12.904} & \emph{---} & \emph{$-8$\%} \\
Round 4e --- \texttt{max-seq-len} 32768 $\to$ 4096         & 13.965 & 0.027 & noise \\
\midrule
\textbf{Stage 11 sub-total} (Rounds 3+4 net)               & \textbf{14.046} & 0.049 & $+$2.38\% vs S8 \\
\midrule
\multicolumn{4}{l}{\emph{2026-05-18 session: four further bit-exact commits pushed to \texttt{pi5-cluster}}} \\
Stage 12 --- remove SW prefetch in matmul (\texttt{64ef787})         & 13.997 & 0.040 & noise \\
Stage 13 --- true single-pass \texttt{silu\_mul\_F32} (\texttt{1af175c}) & 14.270 & --- & $+$1.6\% vs S11 \\
Stage 14 --- \texttt{MAX\_CHUNK\_SIZE} 4K $\to$ 16K (\texttt{f8ed9d2}) & \textbf{14.449} & 0.038 & $+$1.25\% vs S13 \\
Stage 15 --- runtime tweaks persisted via systemd (\texttt{edbc4ce}) & 14.449 & 0.038 & flat \\
\midrule
\textbf{Final, all bit-exact changes (Stages 1--15, mean)} & \textbf{14.449} & \textbf{0.038} & \textbf{$+$5.31\% vs S8, $+$10.81\% vs \#255} \\
\textbf{Best individual run (from same $n$=20)}            & \textbf{14.557} & \emph{---}    & \textbf{$+$11.63\% vs \#255} \\
\bottomrule
\end{tabular}
\end{table}

\subsection{Bottleneck characterisation via ARM PMU}

Profiling with \texttt{perf stat} on both root and worker nodes during sustained inference:

\begin{table}[H]
\centering
\caption{Measured stalls and DRAM traffic (per node, 60s sample, n=4 nodes).}
\small
\begin{tabular}{lrr}
\toprule
\textbf{Metric} & \textbf{Root (rpi-1005)} & \textbf{Worker (rpi-1006)} \\
\midrule
\texttt{stalled-cycles-backend} (\% cycles) & 49.25\% & 47.69\% \\
\texttt{stalled-cycles-frontend} (\% cycles) & 3.83\% & --- \\
DRAM read bandwidth ($l2d\_cache\_refill \times 64$) & 2.74 GB/s & 2.61 GB/s \\
DRAM write bandwidth ($l2d\_cache\_wb \times 64$) & 8.66 GB/s & 8.47 GB/s \\
Total DRAM bandwidth per node & 11.40 GB/s & 11.08 GB/s \\
IPC (insn / cycle) & 1.74 & --- \\
dTLB miss rate & 0.08\% & --- \\
iTLB miss rate & 0.01\% & --- \\
branch-misses & 0.07\% & --- \\
\bottomrule
\end{tabular}
\end{table}

\textbf{The 49\% backend stall, combined with $\sim$11 GB/s per-node DRAM bandwidth ($\approx 85$\% of the practical 13 GB/s 4-thread Triad ceiling for LPDDR4X-4267)}, is the diagnostic smoking gun: the cluster spends roughly half its cycles waiting on memory. Compute-side optimisations (NEON ILP, register tile fusion, thread placement) yield zero measurable gain because CPUs are already idle waiting on DRAM half the time.

The 3.16$\times$ write-to-read ratio (8.66 GB/s writes vs 2.74 GB/s reads) is unusual for streaming-weight inference and points to the multi-pass intermediate buffer pattern in the MoE FFN ($w_3 \to$ silu $\to$ mul produces and writes back temporary buffers between barriers). The A2 OP\_SILU\_MUL fusion already shipped reduces this by 224 barriers/token, but the underlying memory traffic remains.

\section{Round 3 and Round 4 findings}

\subsection{Findings from research agents}

\begin{table}[H]
\centering
\caption{Top findings from 4 background research agents, with applicability to our setup.}
\footnotesize
\begin{tabular}{p{4cm}p{2cm}p{3cm}p{3cm}}
\toprule
\textbf{Technique} & \textbf{Source} & \textbf{Applicable?} & \textbf{Status} \\
\midrule
IRQ suspension (\texttt{irq\_suspend\_timeout}) & LWN 1008399 & No (kernel 6.13+) & Pi5 on 6.12 \\
\texttt{MADV\_HUGEPAGE} & arxiv 2407.13476 et al. & No (THP not compiled in Pi OS) & Tested no-op \\
AF\_XDP zero-copy & kernel docs & No (\texttt{macb} no native XDP) & RFC only \\
io\_uring SQPOLL + DEFER\_TASKRUN & axboe/liburing & Net negative on 4-core & Skipped \\
\texttt{SO\_BUSY\_POLL} + RPS/RFS & LWN 540281 & Already TIER 0 & Applied \\
Threaded NAPI & 5.12+ docs & Tested $-3.9$\% & Reverted \\
SCHED\_FIFO / RR at moderate priority & kernel.org & Considered, deferred & Risk:reward poor \\
NUMA \texttt{fake=4} + \texttt{local} & Linux NUMA docs & Requires reboot & Excluded \\
MPAM cache partitioning & docs.kernel.org & Kernel 6.19+ AND A76 lacks hw & Future kernel \\
Per-task EEVDF \texttt{sched\_runtime} & LWN 969062 & Applied (4ms) & No measurable gain \\
ik\_llama pair-row matmul ILP & PR \#229 follow-up & Tested no-op (OoO already optimal) & Reverted \\
ik\_llama PR \#1485 cross-expert pool & ik\_llama & Already-optimal in dllama & N/A \\
Jumbo frames MTU 9000 & RFC 2861 + macb & bcmgenet EBUSY live link & Documented dead-end \\
Phase B foundation cherry-pick (K=0) & opt-b5b6-tiled-sync & Applied (perf-neutral) & Foundation in repo \\
Phase B \texttt{--tile-sync 4} (no Option C wiring) & opt-b5b6-tiled-sync & Tested $-8\%$ & Reverted \\
\bottomrule
\end{tabular}
\end{table}

\subsection{Phase B foundation cherry-pick}

We cherry-picked four commits from the archived \texttt{opt-b5b6-tiled-sync} branch onto \texttt{pi5-cluster}:

\begin{itemize}[leftmargin=*,nosep]
  \item \texttt{ce7b96e feat(B5+B6): tile-sliced sync behind --tile-sync flag}
  \item \texttt{6ebd00f fix: writeReadMany $\to$ writeMany+readMany}
  \item \texttt{469101e feat(B5+B6): async sync orchestrator (drain thread)}
  \item \texttt{2aba914 feat(B6): matmul\_Q80\_Q40\_F32\_tiled wrapper}
\end{itemize}

With \texttt{--tile-sync 0} (default) the cluster behaviour is byte-identical to the pre-foundation state. We measured 14.092 $\pm$ 0.044 tok/s with K=0, confirming perf-neutrality. With \texttt{--tile-sync 4} active, we measured $-8\%$ regression: the wrapper subdivides syncs into K small writes but \emph{does not yet drive the matmul to emit per-tile signals} (that is the Option~C wiring, $\sim$40 LOC still pending). The TCP small-write penalty without overlap dominates.

\subsection{EEVDF per-task slice tuning}

We added a \texttt{sched\_setattr()} call at the start of \texttt{executorWorkerLoop} requesting a 4ms slice (default 700us under EEVDF in 6.12). We confirmed via \texttt{/proc/$\langle$pid$\rangle$/sched} that workers reported \texttt{se.slice = 4000000}. Measurement: 14.061 $\pm$ 0.048 tok/s, no measurable gain vs the unset case. Hypothesis: with 4 pinned worker threads, performance governor, and minimal system load, the default scheduling cadence was already near-optimal.

\subsection{Sysctls applied (cumulative, persisted)}

\texttt{/etc/sysctl.d/99-dllama-r3.conf} (Round 3):

\begin{lstlisting}
kernel.sched_autogroup_enabled = 0
net.ipv4.tcp_autocorking = 0
net.ipv4.tcp_slow_start_after_idle = 0
net.ipv4.tcp_no_metrics_save = 1
\end{lstlisting}

\texttt{/etc/sysctl.d/99-dllama-r4.conf} (Round 4):

\begin{lstlisting}
net.ipv4.tcp_thin_linear_timeouts = 1
net.core.netdev_max_backlog = 8192
net.core.rmem_max = 33554432
net.core.wmem_max = 33554432
net.ipv4.tcp_rmem = 4096 87380 33554432
net.ipv4.tcp_wmem = 4096 65536 33554432
vm.dirty_bytes = 16777216
vm.dirty_background_bytes = 4194304
\end{lstlisting}

Each individual change produced no measurable gain (within $\pm$0.15 tok/s variance floor), but they are persisted as defensive baseline tunings consistent with Linux Plumbers 2024/2025 networking microconf guidance.

\section{Discussion}

\subsection{Why no further bit-exact improvement?}

The Pi5 LPDDR4X-4267 memory controller has a theoretical peak of 17.07 GB/s and a measured 4-thread Triad ceiling of $\sim$13 GB/s. Our four nodes each sustain $\sim$11 GB/s during inference, or $\approx$85\% of practical peak. The compute path (Q80 $\times$ Q40 matmul with NEON dotprod) is hardware-tuned to the limit; the synchronisation path (TCP all-to-all-broadcast over 1 GbE) is latency-bound at $\sim$125 us per sync round. Together they saturate the system's effective work rate.

\textbf{Conclusion at end of Stage 11}: with the constraints of bit-exact computation, no kernel rebuild, and no quality compromise, the practical ceiling we observed for this hardware running Qwen3-30B-A3B Q40 was $\sim$14.0 tok/s. We achieved 14.046 tok/s reproducibly when the cluster is cold-started.

\textbf{Update at end of Stage 15 (2026-05-18 session)}: the four bit-exact source-level commits described in this paper's updated Table~\ref{tab:trajectory} lift the practical ceiling to \textbf{14.449 tok/s sustained} (best individual run 14.557). Two of the four wins are inversions of long-standing micro-optimisations: Stage~12 removes the \texttt{\_\_builtin\_prefetch} calls from the matmul inner loop (the Cortex-A76's hardware stride detector outperforms the manual hints once the loop is hot), and Stage~13 rewrites the previously two-call SILU+MUL wrapper as a true single-pass NEON kernel that keeps values in registers ($-$33\% memory ops for this op). Stage~14 (the largest single source-level win of the session, \texttt{MAX\_CHUNK\_SIZE} 4 KB $\to$ 16 KB) reduces the \texttt{send()}/\texttt{recv()} syscall count of the all-reduce loop 4-fold. Stage~15 persists the previously runtime-only NIC and softirq tunings (RX ring 4096, RFS, NAPI defer) via a new \texttt{systemd} unit so the operating-point survives reboots.

\subsection{What would unlock further gains?}

\begin{enumerate}[leftmargin=*,nosep]
  \item \textbf{Phase B Option C ($\sim$40 LOC + 60 min soak)}: wire the existing async-sync drain thread to the matmul tile-signal emitter so that the network ships tile $t$ while the matmul computes tile $t+1$. Estimated +5--12\% based on the measured 17\% sync wall and 60\% overlap target. \emph{Foundation is now in pi5-cluster}; only wiring + watchdog remains.
  \item \textbf{Quality compromise}: top-$k$=4 (arxiv:2604.07035) gives +25--40\% on a memory-bandwidth-bound MoE workload, but at the cost of $\sim$1--2 percentage points on reasoning benchmarks. Q3\_K\_S in dllama (would need NEON kernel port) projects +25--35\%. REAP 25\% expert pruning (offline, no LOC change) projects +15--25\%, with $-1.1$\% on math benchmarks.
  \item \textbf{Hardware}: Apple Silicon UMA (M-series) or any platform with $>$30 GB/s effective per-core DRAM. Out of scope.
\end{enumerate}

\subsection{Implications for Pi5 LLM workloads more broadly}

The Pi5 has become a popular substrate for ``edge LLM'' experiments. Our findings suggest:

\begin{itemize}[leftmargin=*,nosep]
  \item For bandwidth-bound workloads (any model larger than $\sim$1 GB of active weights per token per node), the LPDDR4X memory subsystem is the binding constraint. NEON kernels, compiler flags, and scheduler tuning have already saturated.
  \item MoE architectures help disproportionately on Pi5: the same total parameter count with sparse activation maps better to the memory budget than dense.
  \item Network-side tuning (TCP buffers, busy-poll, IRQ affinity) is largely solved in TIER 0; further work on this axis returns within the measurement noise.
  \item Compute/comms overlap (Phase B Option C, FlashOverlap, CommFuse) is the only software-side lever that can hide the memory latency, and it requires dedicated engineering with a safety net (watchdog + soak test).
\end{itemize}

\section{Reproducibility}

All commits are pushed to \href{https://github.com/danielcorrea-hellomatik/distributed-llama}{\texttt{danielcorrea-hellomatik/distributed-llama}} on the \texttt{pi5-cluster} branch:

\begin{table}[H]
\centering
\footnotesize
\begin{tabular}{ll}
\toprule
\textbf{SHA} & \textbf{Change} \\
\midrule
\texttt{162cd83} & docs: SESSION Round 4 final \\
\texttt{2aba914} & feat(B6): matmul\_Q80\_Q40\_F32\_tiled wrapper \\
\texttt{469101e} & feat(B5+B6): async sync orchestrator \\
\texttt{6ebd00f} & fix: writeReadMany $\to$ writeMany+readMany \\
\texttt{ce7b96e} & feat(B5+B6): \texttt{--tile-sync K} flag \\
\texttt{850d1e1} & perf(executor): EEVDF \texttt{sched\_setattr} slice 4ms \\
\texttt{f5bd00c} & docs: SESSION Round 3 (4 background agents) \\
\texttt{33ef826} & docs: SESSION + pair-row matmul ROI 0\% \\
\texttt{c0bcf74} & docs: SESSION EXTENDED initial \\
\texttt{106eab3} & perf(moe): OP\_SILU\_MUL bit-exact ($+$0.5\%) \\
\texttt{d50d30c} & docs(stage9): TIER 0 sysctls ($+$1.5\%) \\
\midrule
\multicolumn{2}{l}{\emph{2026-05-18 session --- Stages 12--15:}} \\
\texttt{64ef787} & perf(matmul): remove SW prefetch (Stage 12) \\
\texttt{1af175c} & perf(silu\_mul): TRUE single-pass NEON kernel (Stage 13, $+$1.6\%) \\
\texttt{f8ed9d2} & perf(network): MAX\_CHUNK\_SIZE 4K $\to$ 16K (Stage 14, $+$1.34\%) \\
\texttt{edbc4ce} & ops(systemd): persist runtime tweaks (Stage 15) \\
\texttt{5092842} & docs(paper): academic-style rewrite + threats + p-values \\
\texttt{7650363} & docs(paper): refresh stale 12.71 references to 14.449 \\
\bottomrule
\end{tabular}
\end{table}

Cluster operational state and full session logs live in \texttt{docs/SESSION-2026-05-17-EXTENDED.md} (488 lines). Per-Pi persisted tunings in \texttt{/etc/sysctl.d/99-dllama.conf} (TIER 0), \texttt{99-dllama-r3.conf} (Round 3), \texttt{99-dllama-r4.conf} (Round 4), and \texttt{/etc/systemd/system/dllama-eth-tune.service} (\texttt{ethtool -K eth0 gro off} at boot).

\section{Conclusion}

We documented a multi-round empirical investigation pushing a 4-node Raspberry Pi 5 cluster running Qwen3-30B-A3B Q40 from its public ceiling of 13.04 tok/s up to \textbf{14.449 tok/s sustained} (best individual run 14.557; $n$=20 paired, 95\% CI $\pm$0.038), a \textbf{+10.81\% improvement} over the state-of-the-art for this exact hardware. Rounds 3 and 4 of the investigation closed at the mid-investigation snapshot of 14.046 tok/s ($+$7.72\%); the subsequent 2026-05-18 session added Stages 12--15 (four further bit-exact source-level commits) that lifted the figure another $+$2.87 percentage points to the final 14.449 tok/s. The journey involved:

\begin{itemize}[leftmargin=*,nosep]
  \item 8 source-level patches to \texttt{distributed-llama} (Stages 1--8).
  \item TIER 0 sysctl bundle (Stage 9), validated at +2.12\% over the patched baseline.
  \item OP\_SILU\_MUL op-fusion (A2), bit-exact +0.5\%.
  \item Phase B async-sync foundation cherry-picked but not activated; Option C wiring remains as future work.
  \item EEVDF per-task slice tuning, two rounds of sysctl bundles, IRQ affinity, per-thread pinning, and madvise hints --- none of which provided measurable additional gain.
  \item Empirical confirmation via ARM PMU that the cluster is DRAM-bandwidth-bound (49\% backend stalls).
  \item \textbf{2026-05-18 Stages 12--15:} four further cumulative bit-exact commits ($+$2.87 percentage points on top of Stage 11). Stage 12 inverts the long-standing software-prefetch micro-optimisation in the matmul inner loop (the Cortex-A76's HW stride detector outperforms the manual hints); Stage 13 rewrites the SILU+MUL fusion as a true single-pass NEON kernel ($-$33\% memory ops for this op); Stage 14 bumps \texttt{MAX\_CHUNK\_SIZE} 4 KB $\to$ 16 KB in the all-reduce loop ($-$4$\times$ \texttt{send()}/\texttt{recv()} syscalls); Stage 15 persists the runtime NIC tuning via \texttt{systemd}. Sustained throughput at the end of Stage 15: \textbf{14.449 tok/s} ($+$10.81\% over public ceiling, $+$5.31\% over Stage 8 baseline).
\end{itemize}

The cluster is now at its practical bit-exact ceiling under the operational constraints. We document the Phase B Option C wiring path as the next sprint candidate ($\sim$40 LOC + watchdog + 60-minute soak test) with an estimated upside of +5--12\%. Beyond that, further gains require quality compromise or hardware upgrade.

\begin{thebibliography}{9}
\small
\bibitem{originalreport} D.\ Correa Villa, ``Distributed LLM Inference on a 4-Node Raspberry Pi 5 Cluster,'' Hellomatik technical report, 16 May 2026. (\texttt{paper/main\_en.tex} in repo)
\bibitem{tadych2024} B.\ Tadych, \emph{distributed-llama}, \url{https://github.com/b4rtaz/distributed-llama}.
\bibitem{dllama255} \emph{Qwen3-30B-A3B on 4$\times$Pi5 8GB --- 13.04 tok/s discussion}, \url{https://github.com/b4rtaz/distributed-llama/discussions/255}.
\bibitem{lwn1008399} J.\ Damato, ``Smarter IRQ suspension,'' LWN 1008399, 2025. \url{https://lwn.net/Articles/1008399/}
\bibitem{lwn969062} ``EEVDF and the per-task slice control,'' LWN 969062. \url{https://lwn.net/Articles/969062/}
\bibitem{lpc2025} Linux Plumbers Conference 2025 networking microconf. \url{https://lpc.events/event/19/sessions/236/}
\bibitem{ikllama} I.\ Kawrakow, \emph{ik\_llama.cpp}, \url{https://github.com/ikawrakow/ik_llama.cpp}.
\bibitem{phaseb} D.\ Correa Villa, ``Phase B plan and wiring,'' \texttt{docs/PHASE-B-PLAN.md} \& \texttt{docs/PHASE-B-WIRING.md} in repo, 2026.
\bibitem{flashoverlap} K.\ Hong, X.\ Li \emph{et al.}, ``FlashOverlap,'' \emph{arXiv:2504.19519}, 2025.
\bibitem{commfuse} ``CommFuse: Hiding tail latency via communication decomposition,'' \emph{arXiv:2604.24013}, 2026.
\bibitem{moemodelpaper} ``Gemma 4, Phi-4, Qwen3 accuracy-efficiency tradeoffs,'' \emph{arXiv:2604.07035}, 2026.
\bibitem{kernel_napi} Linux kernel NAPI documentation. \url{https://docs.kernel.org/networking/napi.html}
\end{thebibliography}

\vfill
\begin{center}
\rule{0.6\textwidth}{0.4pt}\\[0.3cm]
\small\itshape Document generated 17 May 2026 from the operational cluster. Hellomatik · Raspberry Pi Cluster Lab.
\end{center}

\end{document}
